\chapter{Wstęp}
\label{chapter:intro}
\section{Motywacja, cel pracy oraz pytania badawcze}

W ostatnich latach systemy informatyczne są coraz częściej projektowane i wdrażane w środowiskach chmurowych, w których kluczowe znaczenie mają wydajność, skalowalność oraz koszty utrzymania aplikacji. Wśród najczęściej rozważanych podejść architektonicznych architektura monolityczna oraz architektura oparta na mikroserwisach stanowią dwa zasadniczo odmienne sposoby budowy i eksploatacji systemów informatycznych. Architektura monolityczna jest często postrzegana jako prostsza w implementacji, wdrażaniu i utrzymaniu w przypadku mniejszych systemów, natomiast podejście mikroserwisowe jest powszechnie kojarzone z większą modułowością, elastycznością i możliwością skalowania. Korzyści te nie są jednak osiągane bez pewnych kompromisów.\cite{dragoni2017microservicesyesterdaytodaytomorrow}

Rosnąca popularność mikroserwisów sprawiła, że wiele organizacji decyduje się na stosowanie tego stylu architektonicznego, często przy założeniu, że jest on z natury lepszy od architektury monolitycznej, zwłaszcza w środowiskach chmurowych. W praktyce rzeczywiste korzyści wynikające z zastosowania mikroserwisów zależą jednak w dużym stopniu od charakteru aplikacji, rodzaju obciążenia, modelu wdrożenia oraz narzutu operacyjnego. Zwiększona złożoność komunikacji między usługami, orkiestracji, monitorowania oraz zarządzania infrastrukturą może w istotny sposób wpływać zarówno na wydajność systemu, jak i na całkowity koszt jego utrzymania.\cite{waseem2026issues}

Jednocześnie platformy chmurowe oferują mechanizmy takie jak skalowanie wertykalne i horyzontalne, konteneryzacja oraz automatyzacja infrastruktury, które mogą wpływać na efektywność obu stylów architektonicznych. W rezultacie na pytanie, czy architektura monolityczna czy mikroserwisowa sprawdza się lepiej w środowisku chmurowym, nie można udzielić jednoznacznej, uniwersalnej odpowiedzi. Wymaga to przeprowadzenia systematycznego porównania w kontrolowanych warunkach eksperymentalnych, z uwzględnieniem nie tylko podstawowych metryk wydajnościowych, ale również cech skalowalności oraz kosztów utrzymania.\cite{verticalhorizontalscaling}

Mimo że w literaturze przedmiotu można znaleźć liczne opracowania dotyczące wzorców architektonicznych oraz systemów chmurowych, nadal brak jednoznacznego konsensusu co do rzeczywistego wpływu wybranego stylu architektonicznego na wydajność, skalowalność i efektywność kosztową aplikacji wdrażanych w chmurze. Wiele prac koncentruje się wyłącznie na wybranych atrybutach jakościowych, natomiast stosunkowo niewiele przedstawia kompleksowe porównanie implementacji monolitycznej i mikroserwisowej testowanych w porównywalnych warunkach. Z tego względu główną motywacją niniejszej pracy jest zbadanie, w jaki sposób wybór pomiędzy architekturą monolityczną a architekturą opartą na mikroserwisach wpływa na działanie aplikacji wdrożonej w środowisku chmurowym. Analiza została ukierunkowana na trzy kluczowe aspekty: wydajność, skalowalność oraz koszty utrzymania. W celu rozwiązania tego problemu porównane zostaną dwie funkcjonalnie równoważne implementacje tej samej aplikacji: jedna oparta na architekturze monolitycznej, a druga na architekturze mikroserwisowej. Oba warianty zostaną wdrożone w środowisku chmurowym i poddane ocenie z wykorzystaniem zdefiniowanej metodyki benchmarkowej.

Głównym celem pracy jest przeprowadzenie analizy porównawczej wydajności, skalowalności oraz kosztów utrzymania architektury monolitycznej i architektury opartej na mikroserwisach w środowiskach chmurowych. Aby osiągnąć ten cel, przygotowane zostanie środowisko eksperymentalne, zdefiniowane zostaną scenariusze pomiarowe, a uzyskane wyniki zostaną przeanalizowane w celu wskazania mocnych i słabych stron obu podejść.

W pracy sformułowano następujące główne pytanie badawcze:
\begin{quote}
\textit{W jaki sposób wybór pomiędzy architekturą monolityczną a architekturą opartą na mikroserwisach wpływa na wydajność, skalowalność oraz koszty utrzymania aplikacji wdrożonej w środowisku chmurowym?}
\end{quote}

W celu bardziej szczegółowej oceny rozważane są również następujące pytania pomocnicze:
\begin{itemize}
    \item Jak oba style architektoniczne różnią się pod względem przepustowości i czasu odpowiedzi wraz ze wzrostem obciążenia?
    \item Jak skutecznie aplikacje monolityczne i mikroserwisowe wykorzystują skalowanie wertykalne i horyzontalne w środowisku chmurowym?
    \item Jaki wpływ ma wybrany styl architektoniczny na koszty utrzymania systemu?
\end{itemize}

Wyniki przeprowadzonych badań mają bezpośrednie zastosowanie zarówno w analizie naukowej, jak i w praktyce inżynierskiej. W ujęciu naukowym pozwalają na bardziej obiektywną ocenę rzeczywistych różnic pomiędzy architekturą monolityczną a mikroserwisową, w szczególności w kontekście wydajności, skalowalności i kosztów. W ujęciu praktycznym dostarczają konkretnych danych, które mogą być wykorzystane przez architektów oprogramowania i inżynierów do podejmowania świadomych decyzji przy wyborze architektury systemu w środowisku chmurowym.
\section{Struktura pracy}

Poniższa praca składa się z ośmiu rozdziałów. Rozdział~\ref{chapter:intro} wprowadza problem badawczy, cel pracy oraz pytania badawcze. Rozdział~\ref{chapter:archo} przedstawia podstawy teoretyczne architektury monolitycznej i mikroserwisowej, ze szczególnym uwzględnieniem ich wpływu na wydajność, skalowalność, komunikację między usługami oraz koszty utrzymania. W rozdziale~\ref{chapter:slr} zaprezentowano systematyczny przegląd literatury, obejmujący strategię wyszukiwania, kryteria selekcji publikacji, analizę dotychczasowych badań oraz identyfikację luki badawczej. Rozdział~\ref{chapter:implementation} opisuje aplikację benchmarkową, jej wariant monolityczny i mikroserwisowy, wykorzystane technologie oraz sposób przygotowania środowiska eksperymentalnego w Google Cloud Platform. Rozdział~\ref{chapter:experiment_methodology} określa metodykę badań, w tym scenariusze testowe, mierzone metryki, konfiguracje maszyn oraz procedury skalowania wertykalnego i horyzontalnego. W rozdziale~\ref{chapter:results} przedstawiono i przeanalizowano wyniki eksperymentów dotyczące latencji, przepustowości, skuteczności skalowania oraz kosztów środowiska eksperymentalnego. Rozdział~\ref{chapter:dangers} omawia zagrożenia dla rzetelności wyników oraz ograniczenia przyjętego podejścia badawczego. Ostatni rozdział~\ref{chapter:summary}, zawiera odpowiedzi na pytania badawcze, najważniejsze wnioski oraz propozycje dalszych kierunków badań.